\documentclass[10pt,a4paper]{amsart}
\usepackage[margin=19mm]{geometry}
\usepackage{amsmath,amssymb,mathtools}
\usepackage[scr=rsfs,cal=euler]{mathalfa}
\usepackage{xfrac}
\usepackage[unicode,hidelinks,bookmarks=false]{hyperref}
\newcommand{\ie}{i.e.\ }
\newcommand{\cf}{cf.\ }
\newcommand{\resp}{respectively\ }
\newcommand{\Subsection}{\subsection}
\providecommand{\nobreakdash}{\nobreak\hbox{-}}
\setlength{\emergencystretch}{2em}
\multlinegap0pt
\usepackage{bm}
\usepackage{amssymb}
\usepackage{mathtools}
\newtheorem{proposition}{Proposition}
\title{\textbf{Precursors of extreme events and critical transitions}}
\author{Riccardo Consonni, Luca Magri}
\date{Source: arXiv:2604.12869v1 (CC BY 4.0)}
\begin{document}
\maketitle

\noindent\emph{Source and licence.} \textbf{Precursors of extreme events and critical transitions}. Version arXiv:2604.12869v1 (CC BY 4.0), distributed by arXiv under the Creative Commons Attribution 4.0 International licence, http://creativecommons.org/licenses/by/4.0/, as recorded in the arXiv licence metadata for this exact version and shown on the abstract page https://arxiv.org/abs/2604.12869v1. Full licence notice: \texttt{licenses/arxiv-2604.12869-CC-BY-4.0.txt}.\par\smallskip

\noindent\textit{Editorial scope.} Counted unit: the article's proposition 3 (source0.tex lines 258--261). The article's complete original proof of it is delivered with it (source0.tex lines 263--276, one \texttt{proof} environment of 14 lines). No mathematics is written, repaired or re-derived: the statement and the proof are the article's own lines, in the article's own notation, and no retained line is substituted or re-flowed.  Retained uncounted as the source-local context the proof needs: lines 126--128; lines 198--203; lines 352--357.  One declared adaptation: exactly 1 ancillary strings inside the retained spans are omitted (image-inclusion commands and, where the source repeats a label, the repeated label definitions) and no other line differs from the source; the figure environments, with their placement, captions and labels, are kept, and the package records the omitted strings in fidelity receipts bound to the affected bodies.  No retained line contains a citation, so the wrapper carries no bibliography and no bibliography entry.  The retained lines contain the article's own diagram code, which the wrapper typesets with the standard tikz packages; no image file is delivered and the package declares no asset dependency.  Editorial formatting only, in addition: an AMS \texttt{amsart} preamble that restates the article's own declarations and macros used by the retained lines, namely the environments \texttt{proposition} on separate counters, together with the standard packages those lines need.  The retained environments print in this document as Proposition 1 in order of appearance, whereas the article prints the same items with the numbering of its own full document.  The complete native source of the article as distributed in the arXiv e-print for this exact version is bundled unchanged as \texttt{sources/198433/source0.tex}.
\begin{equation} \label{dynamical_system}
\dot{\mathbf{z}} = \mathbf{F}(\mathbf{z}),
\end{equation}

\begin{align} \label{eq_clv projeciton form}
\dot{{\mathbf{u}}} &= \mathbf{J}(t){\mathbf{u}}(t) - \left({{\mathbf{u}}(t)^T \mathbf{J}(t){\mathbf{u}}(t)}\right){\mathbf{u}}(t), \nonumber \\ 
& = \left( \mathbf{I} - \mathbf{u}(t)\mathbf{u}(t)^T \right)\mathbf{J}(t)\mathbf{u} \nonumber \\ 
& = \mathbb{P}(\mathbf{u}(t))\mathbf{J}(t)\mathbf{u}(t) \nonumber \\ 
&=\mathbf{f}_\mathbf{J}(\mathbf{u}(t)),
\end{align}

\begin{figure}
    \centering
    
    \caption{On the left: CLV phase space structure when the eigenvalues of the Jacobian $\mathbf{J}$ of Equation \eqref{dynamical_system} are real and follow the inequality $\mu_1 > \mu_2 > \mu_3$. In the middle: CLV phase space structure when the eigenvalues are real and follow $\mu_1 = \mu_2>\mu_3$. On the right: CLV phase space structure when $\sigma_1 = \sigma_2^*$ are complex conjugate and $\mathrm{Re}(\sigma_1) = \mathrm{Re}(\sigma_2) > \mu_3$. In black, we indicate the fixed points and invariant sets on the foreground hemisphere, in gray we indicate fixed points and invariant sets in the background hemisphere.}
    \label{fig:CLV_sphere}
\end{figure}

\begin{proposition}
\label{proposition 3}
Let $\mathbf{v}$ and ${\mathbf{v}^*}$ be two complex conjugate eigenvectors spanning a real stationary eigenspace $E$. Then the set $M_E = E \bigcap S^{d-1}$ is invariant under the flow associated with $\mathbf{f}_{\mathbf{J}}$ and generates a rotation.
\end{proposition}

\begin{proof}
Let $\sigma = \mu + i\omega$ and ${\sigma}^* = \mu - i\omega$ be the eigenvalues associated with $\mathbf{v}$ and ${\mathbf{v}^*}$, respectively. A real eigenspace $E$ is spanned by a linear combination $a\mathbf{q}+b\mathbf{p}$, where $\mathbf{q} = \mathbf{v} + \mathbf{v}^* = 2\mathrm{Re}(\mathbf{v})$ and $\mathbf{p} = i(\mathbf{v} - {\mathbf{v}}^*)=-2\mathrm{Im}(\mathbf{v})$, and $a,b$ are arbitrary real numbers.
From Equation \eqref{eq_clv projeciton form}

\begin{align}
&\mathbf{f_J}(a\mathbf{q} + b\mathbf{p}) = \nonumber \\ 
&=\mathbf{J} (a\mathbf{q} + b\mathbf{p}) +(a\mathbf{q}+ b\mathbf{p})^T\mathbf{J}(a\mathbf{q} + b\mathbf{p}) (a\mathbf{q} + b\mathbf{p})\nonumber \\ 
&= a(\mu \mathbf{q} + \omega\mathbf{p}) + b(\mu \mathbf{p} - \omega\mathbf{q}) + c(a\mathbf{q} + b\mathbf{p}) , \in E
\end{align}

where $a,b,c$ are real numbers. Since any linear combination of vectors in $E$ is mapped into it, $E$ is invariant under $\mathbf{f_J}$.
Consequently, because the set $M_E = S^{d-1} \bigcap E$ is invariant under the flow map of $\mathbf{f_J}$, and both the sphere $S^{d-1}$ and the eigenspace $E$ are invariant, their intersection must also be invariant.
Hence, $M_E = { \mathbf{u} \in E: ||\mathbf{u}|| = 1 }$ is a circle $S^1$ (an equator) along which the only admissible dynamics are rotations (right panel of Figure \ref{fig:CLV_sphere}).
\end{proof}

\end{document}
